\documentclass[10pt]{article}
\usepackage[utf8]{inputenc}
\usepackage[a4paper,height=20cm,width=13cm]{geometry}
\usepackage{amssymb}
\usepackage{dsfont}
\usepackage{calc}
\usepackage{graphicx}
% \usepackage{pst-node}
\usepackage{fourier}
\usepackage{euscript}
\usepackage{amsmath,amsthm}
\def\phi{\varphi}
\def\theta{\vartheta}
\def\P{\EuScript P}
\def\M{\EuScript M}
\def\B{\EuScript B}
\def\C{\EuScript C}
\def\D{\EuScript D}
\def\U{\EuScript U}
\def\V{\EuScript V}
\def\S{\EuScript S}
\def\sm{\smallsetminus}
\def\niff{\nleftrightarrow}
\def\ZZ{\mathds Z}
\def\NN{\mathds N}
\def\PP{\mathds P}
\def\QQ{\mathds Q}
\def\RR{\mathds R}
\def\<{\langle}
\def\>{\rangle}
\def\E{\exists}
\def\A{\forall}
\def\0{\varnothing}
\def\imp{\rightarrow}
\def\iff{\leftrightarrow}
\def\IMP{\Rightarrow}
\def\IFF{\Leftrightarrow}
\def\range{\textrm{im}}
\def\Mod{\textrm{Mod}}
\def\Aut{\textrm{Aut}}
\def\Th{\textrm{Th}}
\def\acl{\textrm{acl}}
\def\dom{{\rm dom}}
\def\eq{{\rm eq}}
\def\tp{\textrm{tp}}
\def\equivL{\stackrel{\smash{\scalebox{.5}{\rm L}}}{\equiv}}

\def\cnonfork{\mathbin{\raise1.8ex\rlap{\kern0.6ex\rule{0.6ex}{0.1ex}}\rlap{\kern1.1ex\rule{0.1ex}{1.9ex}}\raise-0.3ex\hbox{$\smile$} } }

\newcommand{\labella}[1]{{\sf\footnotesize #1}\hfill}
\renewenvironment{itemize}
  {\begin{list}{$\triangleright$}{%
   \setlength{\parskip}{0mm}
   \setlength{\topsep}{0mm}
   \setlength{\rightmargin}{0mm}
   \setlength{\listparindent}{0mm}
   \setlength{\itemindent}{0mm}
   \setlength{\labelwidth}{3ex}
   \setlength{\itemsep}{0mm}
   \setlength{\parsep}{0mm}
   \setlength{\partopsep}{0mm}
   \setlength{\labelsep}{1ex}
   \setlength{\leftmargin}{\labelwidth+\labelsep}
   \let\makelabel\labella}}{%
   \vspace*{-.3\baselineskip}
  \end{list}}
%\def\ssf#1{\textsf{\small #1}}
\newcounter{ex}
\newenvironment{exercise}{\bigskip\addtocounter{ex}{1}\textbf{Esercizio \theex.\ }}{}
\pagestyle{empty}
\parindent0ex
\parskip1ex
\raggedbottom
\def\nsR{{}^*\!\RR}
\def\ssf#1{\textsf{#1}}
\renewcommand{\baselinestretch}{1.3}


\usepackage{fancyhdr}
\pagestyle{fancy}
\lfoot{Teoria dei Modelli a.a.~2026/27}
\rhead{}
\rfoot{{\small\thepage}}
\cfoot{}

%
% INDICARE NOME E COGNOME DI TUTTI GLI AUTORI
%
\lhead{Nome/i Cognome/i}
%

\begin{document}


\begin{exercise}(Martina Apicella \& Beatrice Calderaro)
Let $M$ be $\omega$-saturated.
Let $A\subseteq M$ be finite.
Prove $M$ realizes every consistent $p(x)\subseteq L(\acl A)$, for $|x|<\omega$.
\end{exercise}

\begin{exercise}(Francesco Dam \& Daniele Lanciotti)
  Prove that if $T$ is $\omega$-categorical then for every $n<\omega$ there is an $m<\omega$ such that $|\acl A|\le m$ for every $A$ such that $|A|\le n$.
\end{exercise}

\begin{exercise}(Filippo Giugie)
Let $p(x)\subseteq L(A)$ and let $\phi(x\,;y)\in L(A)$ be a formula that defines, when restricted to $p(\U^x)$, an equivalence relation with finitely many classes.
Prove that there is a finite equivalence relation definable over $A$ that coincides with $\phi(x\,;y)$ on $p(\U^x)$.
\end{exercise}

\begin{exercise}(Daniele Lanciotti \& Andrea Morello)
Assume that $T$ is strongly minimal.
Let $\phi(x\,;z)\in L(A)$ where $|x|=1$.
For arbitrary $b\in\U^z$, prove that if the orbit of $\phi(\U^x; b)$ over $A$ is finite, then $\phi(\U^x; b)$ is definable over $\acl A$.
\end{exercise}

\end{document}


